\documentclass[9pt]{article}
\usepackage[a4paper,margin=0.25in]{geometry}
\usepackage{amsmath,amssymb,amsfonts,mathtools}
\usepackage{xcolor}
\allowdisplaybreaks
\setlength{\jot}{3pt}

\begin{document}
\section*{The two-$\rho$ terms of the three-$\rho$ equation}
These are the two-$\rho$ terms obtained from the final three-$\rho$ equation (\texttt{ThreeRho\_final.tex}, black and blue lines). They come row by row, in the order of that equation.
\begin{itemize}
\item \textbf{How they are obtained.} Every line is reordered in the charge order written in the three-$\rho$ equation, with $[\rho^a(x),\rho^b(y)]=if^{abc}\rho^c(x)\,\delta^{(2)}(x-y)$:
\begin{align*}
ABC&=S(ABC)+\tfrac14\big(\{[A,B],C\}+\{[A,C],B\}+\{[B,C],A\}\big),\\
A\{B,C\}&=2S(ABC)+\tfrac12\{[A,B],C\}+\tfrac12\{[A,C],B\},\\
\{B,C\}A&=2S(ABC)-\tfrac12\{[A,B],C\}-\tfrac12\{[A,C],B\}.
\end{align*}
The symmetrized parts $S(ABC)$ cancel (three-$\rho$ equation), and one-$\rho$ terms are dropped. The $\delta$-functions are carried out and the colour algebra is done with $U^{bc}=U^{cb}$, $U^{ab}(p)U^{ac}(p)=\delta^{bc}$, $f^{abc}U^{ad}(p)U^{be}(p)U^{cg}(p)=f^{deg}$ and $f^{acd}f^{bcd}=N_c\delta^{ab}$. Terms with the same kernel and the same colour structure are added within each row. The c.c.\ is included.
\item \textbf{Variables.}
\begin{itemize}
\item The measured gluon is at $w$ (amplitude) and $w'$ (c.c.); the unobserved gluon is at $z$.
\item The two remaining charges are $x$ and $x'$. The charge whose emission kernel ends at $w$ is $x$, and the one ending at $w'$ is $x'$. If one charge emits at both $w$ and $w'$, it is called $x'$ and the other one $x$.
\item $\mathcal K(X)=X/X^2$.
\end{itemize}
\item \textbf{Charges.} The two charges form a symmetric product: $\rho^a(p)\,\rho^b(q)\equiv\frac12\{\rho^a(p),\rho^b(q)\}$.
\item \textbf{Coefficients} are real. They multiply $\frac{1}{(2\pi)^3}\frac{g^4}{16\pi^5}\frac1{k^+}\log\frac\vee\Lambda$, so the $i$ of the three-$\rho$ prefactor and the $i$ of the commutator have been multiplied out.
\end{itemize}
There are @@N@@ terms in total. Together with the genuine two-$\rho$ terms of \texttt{Evolution\_two\_rho.tex} (with fixes F5, F6) and the two-$\rho$ lists of \texttt{4rho.tex} (with fixes F1, F2), they give JIMWLK$\otimes$LO exactly. This equation was re-read from this file and checked numerically against that sum.

@@EQ@@
\end{document}
